% Auto-generated by texcv/build.js — do not edit by hand
\documentclass[11pt]{article}
\usepackage[a4paper,top=2.5cm,bottom=2cm,left=2cm,right=2cm]{geometry}
\usepackage{tabularx}
\usepackage{needspace}
\usepackage{parskip}
\usepackage{lmodern}
\usepackage{enumitem}
\usepackage[dvipsnames,table]{xcolor}
\usepackage[T1]{fontenc}
\usepackage[utf8]{inputenc}
\usepackage{hyperref}

\hypersetup{
  colorlinks=true,
  urlcolor=blue,
  linkcolor=black
}

\pagestyle{empty}
\newlength{\cvleftwidth}

\newcommand{\cvsection}[1]{%
  \needspace{4\baselineskip}%
  \vspace{14pt}%
  {\large\textbf{#1}}\\[-4pt]%
  \rule{\linewidth}{0.4pt}\nopagebreak%
  \vspace{3pt}%
}

\begin{document}

\noindent\begin{minipage}[t]{0.5\linewidth}
{\LARGE \textbf{Botond Ortutay}}
\end{minipage}%
\begin{minipage}[t]{0.5\linewidth}
\raggedleft
\textbf{CURRICULUM VITAE}\\
\today
\end{minipage}

\vspace{8pt}

\begin{center}
+358 44 9340057 | \href{mailto:boti.ortutay@gmail.com}{boti.ortutay@gmail.com} | \href{https://www.linkedin.com/in/botond-ortutay/}{LinkedIn} | \href{https://github.com/gitond/}{GitHub}
\end{center}

\cvsection{Profile}
I am currently a student at the University of Turku, although I'll finish my Master's Degree soon. In the meantime I've completed a deep dive into AI, algorithmics, software architecture and all that good stuff. First and foremost I'm a problem solver. If something interests me I find serious joy in building a solution for it.

\cvsection{Education}
\settowidth{\cvleftwidth}{2024--present}\addtolength{\cvleftwidth}{1em}
\begin{tabularx}{\linewidth}{@{}p{\cvleftwidth}X@{}}
  2024--present & \textbf{Master of Science}, University of Turku, Turku, Finland\newline \textit{Major:} Software Engineering\newline \textit{Minor:} Tomorrow's AI\newline \textit{Relevant courses:}\newline \quad \textbullet\ Computer Vision and Sensor Fusion\newline \quad \textbullet\ Textual Data Analysis \\
\end{tabularx}
\vspace{8pt}
\begin{tabularx}{\linewidth}{@{}p{\cvleftwidth}X@{}}
  2020--2024 & \textbf{Bachelor of Science}, University of Turku, Turku, Finland\newline \textit{Major:} Information and Communication Technology\newline \textit{Minor:} Mathematics, Optimization\newline \textit{Relevant courses:}\newline \quad \textbullet\ Web and mobile programming \\
\end{tabularx}
\vspace{8pt}
\begin{tabularx}{\linewidth}{@{}p{\cvleftwidth}X@{}}
  2016--2019 & \textbf{High School}, Tampereen Teknillinen Lukio, Tampere, Finland\newline \textit{Mathematics and IT-focused Matek-line} \\
\end{tabularx}

\cvsection{Publications \& Theses}
\begin{description}
  \item[\textbullet] \textbf{BSc Thesis:} \href{https://www.utupub.fi/server/api/core/bitstreams/5665337e-8e23-4e88-acca-91cfcef3fab7/content}{Polunetsintäalgoritmien esittely ja tehokkuusvertailu}\\
  This thesis examines pathfinding algorithms from a computing-oriented point of view. It is a literature review, so it is mostly based on already published research. It also contains a small set of performance measurements of certain algorithms in a test environment. The thesis introduces the reader to a few pathfinding algorithms, their use cases, and compares their performance.
  \item[\textbullet] \textbf{MSc Thesis:} \href{https://github.com/gitond/msc_thesis_datastructure}{MSc Thesis on Computer Vision integration into Augmented Reality Tracking (work title)} (ongoing)\\
  Ongoing MSc thesis on integrating a computer vision neural network into an AR application to replace "traditional" tracking methods. Current status: Theoretical study (Literature review) completed, working on related software development project.
\end{description}

\cvsection{Work Experience}
\settowidth{\cvleftwidth}{2025--present}\addtolength{\cvleftwidth}{1em}
\begin{tabularx}{\linewidth}{@{}p{\cvleftwidth}X@{}}
  2025--present & \textbf{AI Engineer}, Cloop.io, Turku, Finland\newline The first developer in this company, responsible for designing and building many parts of the current Cloop product, including: RAG (initial version, first vector database, current source scoring system), email system. Active on every part of the stack: Frontend (Chat component, product dashboard), Backend and Database. Relevant technologies: Python (FastAPI, Pydantic), PostgreSQL (incl. pgvector for RAG similarity search), TypeScript. \\
\end{tabularx}
\vspace{8pt}
\begin{tabularx}{\linewidth}{@{}p{\cvleftwidth}X@{}}
  2024--present & \textbf{Project member, Software Engineering \& AI}, VITAR project, Turku, Finland\newline A frozen startup project where we attempted to develop an AI research assistant for biolab workflows. In the beginning I was the only software developer involved in this project. My tasks involved mobile development (Android Studio, Java), software architecture design, and managing AI systems (my first introduction to RAG, Vector databases, N8N and many many more). \\
\end{tabularx}
\vspace{8pt}
\begin{tabularx}{\linewidth}{@{}p{\cvleftwidth}X@{}}
  2023 & \textbf{Project member}, Information Technology Labs, University of Turku, Turku, Finland\newline In this project me and a few other students take part in developing a checkout application to be used with digital learning platforms as a SaaS-service. We used TypeScript and the T3 Stack. My job was to build the backend of the webhook integration. \\
\end{tabularx}
\vspace{8pt}
\begin{tabularx}{\linewidth}{@{}p{\cvleftwidth}X@{}}
  2020 & \textbf{Summer trainee}, HiDucator Ltd, Turku, Finland\newline My main job was to edit educational videos for this e-learning company. I also helped in shooting and producing the raw material for video production. \\
\end{tabularx}

\cvsection{Selected Projects}
\begin{description}
  \item[\textbullet] \textbf{learn3fruit} \href{https://github.com/gitond/learn2fruit}{[repo]} \\\small 2025 to present\quad \colorbox{gray!15}{\strut\small Python} \colorbox{gray!15}{\strut\small TensorFlow} \colorbox{gray!15}{\strut\small TensorFlow.js} \colorbox{gray!15}{\strut\small A-Frame} \colorbox{gray!15}{\strut\small containerization} \colorbox{gray!15}{\strut\small JavaScript}\\This application aims to prove that a Computer Vision Neural Network can be used to power a WebAR/Mobile-AR application. The finished app will guide the user through the preparation of a fruit salad given using a built in recipe (step-by-step format), where the neural network knows where to display the 3D-Augmentations by locating the ingredients \& tools (object detection). This is a part of my MSc Thesis. It currently exists in two git repos, a complete rewrite of what has been developed thus far is in the works.
  \item[\textbullet] \textbf{CV management projects} \href{https://github.com/gitond/cvnew/tree/main}{[repo]} \\\small 2025 to present\quad \colorbox{gray!15}{\strut\small LaTeX} \colorbox{gray!15}{\strut\small Node.js} \colorbox{gray!15}{\strut\small React} \colorbox{gray!15}{\strut\small JavaScript} \colorbox{gray!15}{\strut\small GNU Make}\\A group of git repos that aim to add version control as well as `.tex`, `.pdf` and web renderings of my CV and portfolio.
  \item[\textbullet] \textbf{imageminer/imapp} \textit{(archived)} \href{https://github.com/gitond/imapp}{[repo]} \\\small 2024\quad \colorbox{gray!15}{\strut\small C} \colorbox{gray!15}{\strut\small quirc} \colorbox{gray!15}{\strut\small React Native} \colorbox{gray!15}{\strut\small Android Studio} \colorbox{gray!15}{\strut\small Android NDK}\\Our entry for the 2024 Boost Digital sustainability hackathon. The aim was to come up with a more efficient way to read QR codes on mobile. Our plan was to port over \& compile the quirc C-library for Android using Android's Native Development Kit (NDK), and then embed this in a QR code reading application developed using React Native. This project was never finished, due to a 48h hackathon time limit.
  \item[\textbullet] \textbf{MedMate} \href{https://github.com/gitond/MedMate/tree/main}{[repo]} \\\small 2024\quad \colorbox{gray!15}{\strut\small JavaScript} \colorbox{gray!15}{\strut\small React Native}\\Our concept app for a medication calendar. Developed for the University of Turku course "Lean Digital Business Design". Navigation and UI designs implemented, no real functionality.
  \item[\textbullet] \textbf{graph-tests-2} \href{https://github.com/gitond/graph-tests-2}{[repo]} \\\small 2023--2024\quad \colorbox{gray!15}{\strut\small C++} \colorbox{gray!15}{\strut\small Boost} \colorbox{gray!15}{\strut\small Boost Graph Library} \colorbox{gray!15}{\strut\small R} \colorbox{gray!15}{\strut\small GNU Make}\\Implementations and studying of different pathfinding algorithms. Part of my BSc Thesis Work.
  \item[\textbullet] \textbf{Purgatory} \href{https://v5000a.itch.io/purgatory}{[demo]} \\\small 2023\quad \colorbox{gray!15}{\strut\small Godot}\\My team's submission for the annual GMTK Game Jam in 2023. It was written in the Godot game engine using it's own pythonlike GDscript programming language.
\end{description}

\cvsection{Languages}
\begin{tabularx}{\linewidth}{@{}lX@{}}
  Hungarian & Native \\
  Finnish & Matriculation exam: native (L) \\
  English & Matriculation exam: advanced (M) \\
  Swedish & Matriculation exam: intermediate (M) \\
\end{tabularx}

\cvsection{Programming Languages}
\begin{itemize}[leftmargin=*,noitemsep]
  \item   \textbf{Python} (\textit{Professional}) -- \textbf{libraries}: NumPy, Pandas, PyTorch, TensorFlow, FastAPI, asyncpg, Pydantic; \textbf{package/environment management}: pip, uv, pyenv, venv, jupyter
  \item   \textbf{JavaScript} (\textit{Proficient}) -- \textbf{frameworks}: React, React Native, Vue 3, Node.js; \textbf{libraries}: Transformers.js
  \item   \textbf{C \& C++} (\textit{Familiar}) -- \textbf{libraries}: Boost, Boost Graph Library
  \item   \textbf{Java} (\textit{Familiar})
\end{itemize}

\cvsection{Databases}
\begin{itemize}[leftmargin=*,noitemsep]
  \item   \textbf{SQL} (\textit{Proficient}) -- \textbf{flavors}: PostgreSQL, MariaDB
  \item   \textbf{Vector databases} (\textit{Proficient}) -- \textbf{types}: Qdrant, pgvector for Postgres
  \item   \textbf{MongoDB} (\textit{Familiar})
\end{itemize}

\cvsection{Tools \& Platforms}
\begin{itemize}[leftmargin=*,noitemsep]
  \item \textbf{Linux} (\textit{Proficient}) -- \textbf{distros}: Fedora, Ubuntu, Debian; \textbf{package managers}: dnf, apt
  \item \textbf{Command Line Tools} (\textit{Proficient}) -- \textbf{Examples}: bash, git, make, ssh
  \item \textbf{Containerization tools} (\textit{Proficient}) -- \textbf{environments}: Docker, Docker Compose, Podman
  \item \textbf{AI Tools} (\textit{Proficient}) -- \textbf{libraries}: PyTorch, TensorFlow; \textbf{environments}: Hugging Face; \textbf{LLMs \& code generators}: ChatGPT, Claude Code
  \item \textbf{Figma} (\textit{Familiar})
  \item \textbf{Android Studio} (\textit{Familiar})
\end{itemize}

\end{document}
